% Options for packages loaded elsewhere
\PassOptionsToPackage{unicode}{hyperref}
\PassOptionsToPackage{hyphens}{url}
\PassOptionsToPackage{dvipsnames,svgnames,x11names}{xcolor}
%
\documentclass[
  letterpaper,
  DIV=11,
  numbers=noendperiod]{scrreprt}

\usepackage{amsmath,amssymb}
\usepackage{lmodern}
\usepackage{iftex}
\ifPDFTeX
  \usepackage[T1]{fontenc}
  \usepackage[utf8]{inputenc}
  \usepackage{textcomp} % provide euro and other symbols
\else % if luatex or xetex
  \usepackage{unicode-math}
  \defaultfontfeatures{Scale=MatchLowercase}
  \defaultfontfeatures[\rmfamily]{Ligatures=TeX,Scale=1}
\fi
% Use upquote if available, for straight quotes in verbatim environments
\IfFileExists{upquote.sty}{\usepackage{upquote}}{}
\IfFileExists{microtype.sty}{% use microtype if available
  \usepackage[]{microtype}
  \UseMicrotypeSet[protrusion]{basicmath} % disable protrusion for tt fonts
}{}
\makeatletter
\@ifundefined{KOMAClassName}{% if non-KOMA class
  \IfFileExists{parskip.sty}{%
    \usepackage{parskip}
  }{% else
    \setlength{\parindent}{0pt}
    \setlength{\parskip}{6pt plus 2pt minus 1pt}}
}{% if KOMA class
  \KOMAoptions{parskip=half}}
\makeatother
\usepackage{xcolor}
\setlength{\emergencystretch}{3em} % prevent overfull lines
\setcounter{secnumdepth}{5}
% Make \paragraph and \subparagraph free-standing
\ifx\paragraph\undefined\else
  \let\oldparagraph\paragraph
  \renewcommand{\paragraph}[1]{\oldparagraph{#1}\mbox{}}
\fi
\ifx\subparagraph\undefined\else
  \let\oldsubparagraph\subparagraph
  \renewcommand{\subparagraph}[1]{\oldsubparagraph{#1}\mbox{}}
\fi

\usepackage{color}
\usepackage{fancyvrb}
\newcommand{\VerbBar}{|}
\newcommand{\VERB}{\Verb[commandchars=\\\{\}]}
\DefineVerbatimEnvironment{Highlighting}{Verbatim}{commandchars=\\\{\}}
% Add ',fontsize=\small' for more characters per line
\usepackage{framed}
\definecolor{shadecolor}{RGB}{241,243,245}
\newenvironment{Shaded}{\begin{snugshade}}{\end{snugshade}}
\newcommand{\AlertTok}[1]{\textcolor[rgb]{0.68,0.00,0.00}{#1}}
\newcommand{\AnnotationTok}[1]{\textcolor[rgb]{0.37,0.37,0.37}{#1}}
\newcommand{\AttributeTok}[1]{\textcolor[rgb]{0.40,0.45,0.13}{#1}}
\newcommand{\BaseNTok}[1]{\textcolor[rgb]{0.68,0.00,0.00}{#1}}
\newcommand{\BuiltInTok}[1]{\textcolor[rgb]{0.00,0.23,0.31}{#1}}
\newcommand{\CharTok}[1]{\textcolor[rgb]{0.13,0.47,0.30}{#1}}
\newcommand{\CommentTok}[1]{\textcolor[rgb]{0.37,0.37,0.37}{#1}}
\newcommand{\CommentVarTok}[1]{\textcolor[rgb]{0.37,0.37,0.37}{\textit{#1}}}
\newcommand{\ConstantTok}[1]{\textcolor[rgb]{0.56,0.35,0.01}{#1}}
\newcommand{\ControlFlowTok}[1]{\textcolor[rgb]{0.00,0.23,0.31}{#1}}
\newcommand{\DataTypeTok}[1]{\textcolor[rgb]{0.68,0.00,0.00}{#1}}
\newcommand{\DecValTok}[1]{\textcolor[rgb]{0.68,0.00,0.00}{#1}}
\newcommand{\DocumentationTok}[1]{\textcolor[rgb]{0.37,0.37,0.37}{\textit{#1}}}
\newcommand{\ErrorTok}[1]{\textcolor[rgb]{0.68,0.00,0.00}{#1}}
\newcommand{\ExtensionTok}[1]{\textcolor[rgb]{0.00,0.23,0.31}{#1}}
\newcommand{\FloatTok}[1]{\textcolor[rgb]{0.68,0.00,0.00}{#1}}
\newcommand{\FunctionTok}[1]{\textcolor[rgb]{0.28,0.35,0.67}{#1}}
\newcommand{\ImportTok}[1]{\textcolor[rgb]{0.00,0.46,0.62}{#1}}
\newcommand{\InformationTok}[1]{\textcolor[rgb]{0.37,0.37,0.37}{#1}}
\newcommand{\KeywordTok}[1]{\textcolor[rgb]{0.00,0.23,0.31}{#1}}
\newcommand{\NormalTok}[1]{\textcolor[rgb]{0.00,0.23,0.31}{#1}}
\newcommand{\OperatorTok}[1]{\textcolor[rgb]{0.37,0.37,0.37}{#1}}
\newcommand{\OtherTok}[1]{\textcolor[rgb]{0.00,0.23,0.31}{#1}}
\newcommand{\PreprocessorTok}[1]{\textcolor[rgb]{0.68,0.00,0.00}{#1}}
\newcommand{\RegionMarkerTok}[1]{\textcolor[rgb]{0.00,0.23,0.31}{#1}}
\newcommand{\SpecialCharTok}[1]{\textcolor[rgb]{0.37,0.37,0.37}{#1}}
\newcommand{\SpecialStringTok}[1]{\textcolor[rgb]{0.13,0.47,0.30}{#1}}
\newcommand{\StringTok}[1]{\textcolor[rgb]{0.13,0.47,0.30}{#1}}
\newcommand{\VariableTok}[1]{\textcolor[rgb]{0.07,0.07,0.07}{#1}}
\newcommand{\VerbatimStringTok}[1]{\textcolor[rgb]{0.13,0.47,0.30}{#1}}
\newcommand{\WarningTok}[1]{\textcolor[rgb]{0.37,0.37,0.37}{\textit{#1}}}

\providecommand{\tightlist}{%
  \setlength{\itemsep}{0pt}\setlength{\parskip}{0pt}}\usepackage{longtable,booktabs,array}
\usepackage{calc} % for calculating minipage widths
% Correct order of tables after \paragraph or \subparagraph
\usepackage{etoolbox}
\makeatletter
\patchcmd\longtable{\par}{\if@noskipsec\mbox{}\fi\par}{}{}
\makeatother
% Allow footnotes in longtable head/foot
\IfFileExists{footnotehyper.sty}{\usepackage{footnotehyper}}{\usepackage{footnote}}
\makesavenoteenv{longtable}
\usepackage{graphicx}
\makeatletter
\def\maxwidth{\ifdim\Gin@nat@width>\linewidth\linewidth\else\Gin@nat@width\fi}
\def\maxheight{\ifdim\Gin@nat@height>\textheight\textheight\else\Gin@nat@height\fi}
\makeatother
% Scale images if necessary, so that they will not overflow the page
% margins by default, and it is still possible to overwrite the defaults
% using explicit options in \includegraphics[width, height, ...]{}
\setkeys{Gin}{width=\maxwidth,height=\maxheight,keepaspectratio}
% Set default figure placement to htbp
\makeatletter
\def\fps@figure{htbp}
\makeatother
\newlength{\cslhangindent}
\setlength{\cslhangindent}{1.5em}
\newlength{\csllabelwidth}
\setlength{\csllabelwidth}{3em}
\newlength{\cslentryspacingunit} % times entry-spacing
\setlength{\cslentryspacingunit}{\parskip}
\newenvironment{CSLReferences}[2] % #1 hanging-ident, #2 entry spacing
 {% don't indent paragraphs
  \setlength{\parindent}{0pt}
  % turn on hanging indent if param 1 is 1
  \ifodd #1
  \let\oldpar\par
  \def\par{\hangindent=\cslhangindent\oldpar}
  \fi
  % set entry spacing
  \setlength{\parskip}{#2\cslentryspacingunit}
 }%
 {}
\usepackage{calc}
\newcommand{\CSLBlock}[1]{#1\hfill\break}
\newcommand{\CSLLeftMargin}[1]{\parbox[t]{\csllabelwidth}{#1}}
\newcommand{\CSLRightInline}[1]{\parbox[t]{\linewidth - \csllabelwidth}{#1}\break}
\newcommand{\CSLIndent}[1]{\hspace{\cslhangindent}#1}

\KOMAoption{captions}{tableheading}
\makeatletter
\makeatother
\makeatletter
\@ifpackageloaded{bookmark}{}{\usepackage{bookmark}}
\makeatother
\makeatletter
\@ifpackageloaded{caption}{}{\usepackage{caption}}
\AtBeginDocument{%
\ifdefined\contentsname
  \renewcommand*\contentsname{Table of contents}
\else
  \newcommand\contentsname{Table of contents}
\fi
\ifdefined\listfigurename
  \renewcommand*\listfigurename{List of Figures}
\else
  \newcommand\listfigurename{List of Figures}
\fi
\ifdefined\listtablename
  \renewcommand*\listtablename{List of Tables}
\else
  \newcommand\listtablename{List of Tables}
\fi
\ifdefined\figurename
  \renewcommand*\figurename{Figure}
\else
  \newcommand\figurename{Figure}
\fi
\ifdefined\tablename
  \renewcommand*\tablename{Table}
\else
  \newcommand\tablename{Table}
\fi
}
\@ifpackageloaded{float}{}{\usepackage{float}}
\floatstyle{ruled}
\@ifundefined{c@chapter}{\newfloat{codelisting}{h}{lop}}{\newfloat{codelisting}{h}{lop}[chapter]}
\floatname{codelisting}{Listing}
\newcommand*\listoflistings{\listof{codelisting}{List of Listings}}
\makeatother
\makeatletter
\@ifpackageloaded{caption}{}{\usepackage{caption}}
\@ifpackageloaded{subcaption}{}{\usepackage{subcaption}}
\makeatother
\makeatletter
\@ifpackageloaded{tcolorbox}{}{\usepackage[many]{tcolorbox}}
\makeatother
\makeatletter
\@ifundefined{shadecolor}{\definecolor{shadecolor}{rgb}{.97, .97, .97}}
\makeatother
\makeatletter
\makeatother
\ifLuaTeX
  \usepackage{selnolig}  % disable illegal ligatures
\fi
\IfFileExists{bookmark.sty}{\usepackage{bookmark}}{\usepackage{hyperref}}
\IfFileExists{xurl.sty}{\usepackage{xurl}}{} % add URL line breaks if available
\urlstyle{same} % disable monospaced font for URLs
\hypersetup{
  pdftitle={Tampa Bay Open Science Workshop},
  pdfauthor={Dr.~Chris Anastasiou and Dr.~Marcus Beck},
  colorlinks=true,
  linkcolor={blue},
  filecolor={Maroon},
  citecolor={Blue},
  urlcolor={Blue},
  pdfcreator={LaTeX via pandoc}}

\title{Tampa Bay Open Science Workshop}
\author{Dr.~Chris Anastasiou and Dr.~Marcus Beck}
\date{2022-09-19T10:13:45-04:00}

\begin{document}
\maketitle
\ifdefined\Shaded\renewenvironment{Shaded}{\begin{tcolorbox}[frame hidden, boxrule=0pt, borderline west={3pt}{0pt}{shadecolor}, enhanced, breakable, interior hidden, sharp corners]}{\end{tcolorbox}}\fi

\renewcommand*\contentsname{Table of contents}
{
\hypersetup{linkcolor=}
\setcounter{tocdepth}{2}
\tableofcontents
}
\bookmarksetup{startatroot}

\hypertarget{course-synopsis}{%
\chapter*{Course synopsis}\label{course-synopsis}}
\addcontentsline{toc}{chapter}{Course synopsis}

\includegraphics{./img/2023 Open Science Workshop.png}

Welcome to the Tampa Bay open science workshop: core concepts for
impactful research and resource management. Open science (OS) has been
advocated as an effective approach to create reproducible, transparent,
and actionable research products. However, widespread adoption among the
coastal and estuarine research community has not occurred despite its
perceived benefits. In the face of major challenges like global warming
and sea level rise, the collaborative framework provided by OS is needed
now more than ever. This workshop will cover a half-day of material
introducing participants to core concepts of OS. The target audience
includes anyone interested in applying OS in their own workflows as part
of a larger research and resource management team.

By the end of this workshop, you should have a solid understanding of
fundamental concepts in open science and how they can be applied to help
bridge the research-management divide. You will also have the skills to
understand how collaborative open science tools can be used to increase
efficiency and transparency, understand fundamental best practices for
working with data to facilitate openness, and be able to apply these
lessons within your own teams by effectively addressing barriers to
adoption.

\hypertarget{agenda}{%
\section*{Agenda}\label{agenda}}
\addcontentsline{toc}{section}{Agenda}

\begin{enumerate}
\def\labelenumi{\arabic{enumi}.}
\tightlist
\item
  \protect\hyperlink{basics}{The basics of open science} (\textbf{30
  min})
\item
  \protect\hyperlink{collaborate}{Open science for collaboration}
  (\textbf{90 min})
\item
  \protect\hyperlink{impact}{Open science for impactful products}
  (\textbf{90 min})
\item
  \protect\hyperlink{implement}{Lowering barriers to inclusion and
  addressing key critiques} (\textbf{30 min})
\end{enumerate}

Each module uses a set of common icons to orient you to specific tasks
or experiences during this workshop. These include the following:

\includegraphics[width=1em,height=1em]{./index_files/figure-pdf/fa-icon-20d474448f872ee43905e611a2502347.pdf}
Exercise and discussion

\includegraphics[width=1.25em,height=1em]{./index_files/figure-pdf/fa-icon-5a0e5cd7bfbf585f8546c1ff4cebae73.pdf}
Watch and learn

\includegraphics[width=0.75em,height=1em]{./index_files/figure-pdf/fa-icon-8447cf12acba1bc263badfefa79f6dd7.pdf}
Description of a collaborative tool

\includegraphics[width=1em,height=1em]{./index_files/figure-pdf/fa-icon-e4ee65476be467d7be8a1ae9cb02ffda.pdf}
Pros of a collaborative tool or solution to an open science challenge

\includegraphics[width=1em,height=1em]{./index_files/figure-pdf/fa-icon-c9ddf618df48250f2c0ed0cb59d4f20d.pdf}
Cons of a collaborative tool

\includegraphics[width=1em,height=1em]{./index_files/figure-pdf/fa-icon-e38cc40f3c6469ef5f44a7764a5c5910.pdf}
Challenge to overcome for open science

\hypertarget{instructors}{%
\section*{Instructors}\label{instructors}}
\addcontentsline{toc}{section}{Instructors}

\href{mailto:chris.anastasiou@swfwmd.state.fl.us}{Dr.~Chris Anastasiou}
is a Chief Water Quality Scientist and the Seagrass Mapping Program Lead
for the Southwest Florida Water Management District. He holds a PhD in
Marine Science from the University of South Florida and has been working
in and around the springs of Florida for more than 25 years.

\href{mailto:mbeck@tbep.org}{Dr.~Marcus Beck} is the Program Scientist
for the Tampa Bay Estuary Program and is developing data analysis and
visualization methods for Bay health indicators. He received his PhD in
Conservation Biology from the University of Minnesota in 2013. Marcus
has experience researching environmental indicators and developing open
science products to support environmental decision-making. Marcus is
also an open source software and dashboard developer to facilitate
science application.

\begin{center}\rule{0.5\linewidth}{0.5pt}\end{center}

~~This book is licensed under a Creative Commons Attribution 4.0
International License.

\bookmarksetup{startatroot}

\hypertarget{basics}{%
\chapter{Basics of open science}\label{basics}}

\hypertarget{lesson-outline}{%
\section{Lesson Outline}\label{lesson-outline}}

\begin{itemize}
\tightlist
\item
  \protect\hyperlink{goals-and-motivation-3}{Goals and motivation}
\item
  \protect\hyperlink{why-open-science}{Why open science?}
\item
  \protect\hyperlink{learning-and-speaking-the-language-of-open-science}{Learning
  and speaking the language of open science}
\item
  \protect\hyperlink{fair}{The FAIR principles}
\item
  \protect\hyperlink{schools}{Schools of thought}
\end{itemize}

\hypertarget{goals-and-motivation}{%
\section{Goals and motivation}\label{goals-and-motivation}}

This is the first module in our four hour workshop on open science. This
module describes the need for open science, how it can improve research
applications, and exposes you to common ideas and terminology that we'll
be using throughout the day. Consider this your 30,000 foot view of open
science. Our later modules will provide more detail on specific topics
in open science that you can use for continued learning.

\begin{itemize}
\tightlist
\item
  \textbf{Goal}: get comfortable with key ideas and concepts for
  understanding open science
\item
  \textbf{Motivation}: This is the first step in your open science
  journey!
\end{itemize}

\hypertarget{why-open-science}{%
\section{Why open science?}\label{why-open-science}}

Let's start with revisiting the scientific process. I'm sure this looks
familiar to all of you. This is geared towards an applied research
question.

\includegraphics{./img/scipro.png}

Our basic scientific approach to discovery is motivated by a question or
research goal, developing a hypothesis for the question, collecting data
based on the hypothesis, developing a tool that can be used for
decision-making, and summarizing the results in a conventional format.

In a little more detail, your workflow may look something like this.

\includegraphics{./img/appscipro.png}

Many scientists, especially early career researchers (my past self
included), may assume that this is sufficient to affect change. We write
the report, send it out into the world, and move on to the next project.
This is a common mentality:

\begin{quote}
``This 500-page report will answer all of their questions!''
\end{quote}

From the other side, such as the manager or policy-maker, the report may
be received like this:

\begin{quote}
``This 500-page report answers none of my questions!''
\end{quote}

It's dense, inaccessible, and there are probably questions about the
underlying data and methods used to achieve the results. More
importantly, it doesn't present the information in an easily digestible
format to quickly make the right decision. Sometimes, if you think
you're doing applied science, it may just be \emph{implied} science that
falls short of application.

\textbf{Why is this conventional approach to science ineffective at
seeding change?}

As I'm sure we're aware, the environmental management community is often
siloed with each branch doing their own thing and speaking their own
language. Between the research (typically academic) and management
community, we call this the research-management divide.

\includegraphics{./img/divide2.png}

A distinct gap exists between how scientific products are developed and
how they can be used to meet management needs. This is often the result
of communication barriers, irreproducible results, information loss with
poor documentation, inaccessible data, and opaque workflows known only
to the analyst.

These barriers can occur at any stage of the research process. This
compelling graphic from Michener et al. (1997) describes the atrophy of
information in a closed approach to creating science.

\includegraphics{./img/michener.png}

The last part is especially morbid. Sometimes, this is called the ``bus
factor''. What would happen to your important work and life achievements
if you were hit by a bus? Would others be able to pick it up? Research
products with a high bus factor are at risk of being lost if critical
team members are no longer available. This is a very real problem for
continuity of science.

\includegraphics{./img/overlyhonest.png}

So how do we make changes to our workflows to ensure we can achieve
truly applied science using open tools and philosophies?

\hypertarget{learning-and-speaking-the-language-of-open-science}{%
\section{Learning and speaking the language of open
science}\label{learning-and-speaking-the-language-of-open-science}}

The tools and broader philosophy behind open science can help us bridge
the research-management divide. It involves a fundamental shift in how
we approach the scientific process, both for your own internal workflows
and how you can engage others in the process. By others, we mean not
just researchers, but specifically those that need the information to
make informed decisions.

Before we present a formal definition, let's describe a modification of
the conventional workflow that includes an open process to discovery and
implementation (Beck et al. 2020; modified from Hampton et al. 2015).

\includegraphics{./img/open.png}

This workflow is similar to the original scientific method, but the
technical components are open to managers and stakeholders, we're
treating data differently by using metadata and archiving, we're
creating summary documents that include source code with text, and we're
producing decision-support tools to meet the needs outside of the
research community. More importantly, the process is \emph{iterative}
and not static.

Throughout this workshop, we'll learn about some open science tools that
can be used in this generalized workflow to achieve better science in
less time (Lowndes et al. 2017).

Now let's settle on a definition for open science (from Open Knowledge
International, \url{http://opendefinition.org/},
\url{https://creativecommons.org/}):

\begin{quote}
``The practice of science in such a way that others can
\textbf{collaborate} and \textbf{contribute}, where research data, lab
notes and other research processes are \textbf{freely available}, under
terms that enable \textbf{reuse}, \textbf{redistribution} and
\textbf{reproduction} of the research and its underlying data and
methods.''
\end{quote}

Key words from this definition are highlighted in bold. There are very
specific tools in the open science toolbox that speak to each of these
key words. We'll cover some of these later.

From these key words, we can breakdown this definition into key
principles.

\begin{enumerate}
\def\labelenumi{\arabic{enumi}.}
\item
  \includegraphics[width=1em,height=1em]{./basics_files/figure-pdf/fa-icon-027224388e66831157840c588dc43741.pdf}
  \textbf{Open data}

  \begin{itemize}
  \tightlist
  \item
    Public availability of data
  \item
    Reusability and transparent workflows
  \item
    Data provenance and metadata
  \end{itemize}
\item
  \includegraphics[width=1.25em,height=1em]{./basics_files/figure-pdf/fa-icon-f702081e1a96180027748dfa6dd275a1.pdf}
  \textbf{Open process}

  \begin{itemize}
  \tightlist
  \item
    Iterative methods using reproducible workflows
  \item
    Collaboration with colleagues using web-based tools
  \item
    Leveraging external, open-source applications
  \end{itemize}
\item
  \includegraphics[width=0.88em,height=1em]{./basics_files/figure-pdf/fa-icon-ba6f5f2818a1e4c57b85400a286bb74c.pdf}
  \textbf{Open products}

  \begin{itemize}
  \tightlist
  \item
    Interactive web products for communication
  \item
    Dynamic documents with source code
  \item
    Integration with external networks for discoverability
  \end{itemize}
\end{enumerate}

You'll notice that web-based tools and open science are often discussed
at the same time. Science existed before the internet and open science
often focuses on how the two can leverage and support one another
despite the latter being a relatively new addition to society. People
often describe web-based tools as synonymous with open science.

\hypertarget{fair}{%
\section{The FAIR principles}\label{fair}}

Advocates of open science also describe the use of the FAIR principles
(Wilkinson et al. 2016) as a vehicle for achieving the former. It's
important to understand what they mean so that you can be fluent in
both. The FAIR acronym is described as follows:

\begin{itemize}
\tightlist
\item
  \textbf{F}indable: The data have a globally unique and persistent
  identifier, including use of ``rich'' metadata.
\item
  \textbf{A}ccessible: Once found, the data can be retrieved using
  standardized communications protocols that are open, free, and
  universally implementable.
\item
  \textbf{I}nteroperable: The ability of data or tools from
  non-cooperating resources to integrate or work together with minimal
  effort.
\item
  \textbf{R}eusable: If the above are achieved, the data and metadata
  are described in a way that they can be replicated and/or combined in
  different settings.
\end{itemize}

Simply, what this means is: 1) each dataset has a name that doesn't
change and can be found with minimal effort using that name, 2) once
it's found, you can actually get your hands on it (e.g., not behind a
paywall), 3) once you have it, you can use readily available tools to
work with the data (e.g., not using proprietary software), and 4) you
can actually apply the data for your own needs because it has sufficient
context, including its reproduction, given that the first three
principles are met.

For our purposes, think of these ideas as general guidelines you can ask
yourself when doing science. If you find that your work does not achieve
FAIR status, then you're probably not being as open as you could be.
We'll of course provide some tools to help you be FAIR and open.

\hypertarget{schools}{%
\section{Schools of thought}\label{schools}}

Finally, it's useful to make a distinction of how different people may
talk about open science. This can help you better navigate conversations
and become an advocate for open science in your own right.

A useful paradigm is provided by Fecher and Friesike (2014) to describe
open science as five distinct schools of thought:

\includegraphics{./img/fiveschools.PNG}

These are of course only conceptual boxes and there's considerable
overlap across all schools when open science is used in practice. For
our purposes, we'll be talking about ideas and tools from the pragmatic,
infrastructure, and democratic schools of the thought. The end goal is
to provide you with the means to create more efficient and impactful
science that can more readily be used by others in a collaborative
setting.

\includegraphics[width=1em,height=1em]{./basics_files/figure-pdf/fa-icon-20d474448f872ee43905e611a2502347.pdf}
Exercise and discussion

How do you currently define open science? Has any of the above changed
your understanding of what open science means for research and/or
resource management?

\bookmarksetup{startatroot}

\hypertarget{collaborate}{%
\chapter{Open science for collaboration}\label{collaborate}}

\hypertarget{lesson-outline-1}{%
\section{Lesson Outline}\label{lesson-outline-1}}

\begin{itemize}
\tightlist
\item
  \protect\hyperlink{goals-and-motivation-3}{Goals and motivation}
\item
  \protect\hyperlink{essential-elements-of-collaboration}{Essential
  elements of collaboration}
\item
  \protect\hyperlink{tools-for-collaboration}{Tools for collaboration}
\item
  \protect\hyperlink{balancing-and-managing-expectations}{Balancing and
  managing expectations}
\end{itemize}

\hypertarget{goals-and-motivation-1}{%
\section{Goals and motivation}\label{goals-and-motivation-1}}

This is the second module in our four hour workshop on open science.
This module will explore some open science tools to help you and your
team become better collaborators and to better engage your science with
external partners. We'll introduce some essential elements of
collaboration and discuss some readily available tools for doing so.
We'll close with a discussion on balancing and managing expectations for
collaboration in the rapidly evolving online arena.

\begin{itemize}
\tightlist
\item
  \textbf{Goal}: understand methods of collaboration and the pros/cons
  of various tools
\item
  \textbf{Motivation}: start building the tools for your open science
  toolbox
\end{itemize}

\hypertarget{essential-elements-of-collaboration}{%
\section{Essential elements of
collaboration}\label{essential-elements-of-collaboration}}

We start our deep dive into open science by focusing on collaboration as
the single most important activity that can be enhanced through
transparent, efficient, and reproducible tools. Having effective tools
to work together is a critical theme of many open science practices.
There are many tools in the toolbox and we need to introduce some core
concepts before we demonstrate how to implement them in practice.

\hypertarget{workflow-management}{%
\subsection{Workflow management}\label{workflow-management}}

How do you organize your work each day? How do you make sure pressing
deadlines are met on schedule? How do you plan for short-term and
long-term goals? Do you have a five-year, ten-year, or longer career
plan?

Work to achieve goals cannot be accomplished without a systematic
approach to organizing tasks. Chances are, we each have our own system
that works for us that was developed through trial and error. Although
everyone has familiar workflows, they are often idiosyncratic and deeply
entrenched by habit. That can be in direct conflict with collaboration
when we try to mesh internal workflows with those of others.

Does this look familiar?

\begin{figure}

{\centering \includegraphics[width=0.5\textwidth,height=\textheight]{./img/old_files.png}

}

\end{figure}

Although the above comic from \href{https://xkcd.com/}{xkcd} speaks
directly to file management, it hints at a broader problem of personal
information management that can seriously complicate working with
others. I'm sure we've all struggled to find that one file for that one
project from a vague recollection of seeing it a few months ago.

Collaborative work can be facilitated through workflow management that
helps you break out of old habits. We'll introduce some specific
internet-based tools below to facilitate workflows either for yourself
or, better yet, working with others. These can help propel you towards
open science.

Here, we introduce the
\href{https://en.wikipedia.org/wiki/Kanban_(development)}{Kanban}
approach to workflow management. The idea is simple. Create a
task-oriented workflow using a card management system organized by
progress. It looks something like this:

\begin{figure}

{\centering \includegraphics[width=0.6\textwidth,height=\textheight]{./img/kanban.jpg}

}

\end{figure}

As shown, this approach can work as a literal, physical board or as one
used digitally through a web browser or other software. Every Kanban
board has the following elements that allow you to work in a more
informed manner:

\begin{itemize}
\tightlist
\item
  Provides a ``big picture'' of progress
\item
  Organizes progress by discrete steps
\item
  Establishes cards as specific tasks
\end{itemize}

Many of the open science tools we describe below use this system. It is
a generalizable format that works in different settings, whether it be
general project management or something more formal like software
development.

\includegraphics[width=1.25em,height=1em]{./collaborate_files/figure-pdf/fa-icon-5a0e5cd7bfbf585f8546c1ff4cebae73.pdf}
Watch and learn

Workflow management in the real world - using GitHub to collaborate.

\hypertarget{version-control}{%
\subsection{Version control}\label{version-control}}

A specific problem for workflow management that can be solved by open
science tools is file management. Workflows can greatly be enhanced by
tools that use strict guidelines for tracking changes and allowing a
complete view of the evolution of a project. This is where version
control comes in.

I'm sure many of you have fallen into this trap:

\begin{figure}

{\centering \includegraphics[width=0.5\textwidth,height=\textheight]{./img/filehell.png}

}

\end{figure}

Version control is a way to track the development history of a project.
It serves the joint purposes of:

\begin{enumerate}
\def\labelenumi{\arabic{enumi}.}
\tightlist
\item
  Formally documenting the changes that have been made to code or
  software
\item
  Making sure that the development history is permanent
\item
  Providing a system for collaborating across platforms (with friends!)
\end{enumerate}

It's more than saving files. Documenting changes with a set of commands
that follow strict rules provides a transparent record for yourself and
others, and establishing permanency ensures that any of the changes that
are made can be vetted and accessed as needed. Think of it as an
insurance plan for your project.

By far, the most widely use software for version control is
\href{https://git-scm.com/}{Git}. Although we do not cover the specifics
of this software, it's useful to understand the purpose and what it can
do in making your work more open and impactful. Git is integrated with
many popular open source development platforms, such as
\href{https://support.rstudio.com/hc/en-us/articles/200532077-Version-Control-with-Git-and-SVN}{RStudio}.

Many people often confuse Git with
\href{https://www.github.com}{GitHub}. GitHub is an online platform for
working collaboratively through Git AND it allows you to be open with
your work. We'll provide some examples below of how this can be done.
Importantly, you do not need to be an expert in Git to be able to use
GitHub. This speaks volumes for how team efficiency can be improved
through better collaboration.

\includegraphics[width=1.25em,height=1em]{./collaborate_files/figure-pdf/fa-icon-5a0e5cd7bfbf585f8546c1ff4cebae73.pdf}
Watch and learn

Now we'll demonstrate how to setup a version control project with
RStudio, Git, and GitHub.

\hypertarget{data-repositories}{%
\subsection{Data repositories}\label{data-repositories}}

How data are treated as living, dynamic pieces of information is
critical to the whole ethos of open science. This is especially true
when the \protect\hyperlink{fair}{FAIR} principles are invoked. Data
should not live on your hard drive as something only known to yourself.

Although we will not cover data repositories in depth, it's important to
recognize the critical role that data archiving and metadata have in
open science. How many times have you thought ``wow, it would be great
if I could have the data from this paper!'' Making data open is a great
way to propel science through better collaboration.

The ease of getting a dataset online depends on where you want to put
the data. In all cases, your dataset should be tidy and accompanied by
metadata. For simple solutions, such as FTP hosting or putting a dataset
on Google Drive, all you need to do is upload the data by hand. However,
this doesn't necessarily make it findable and the permanency is
uncertain.

The absolute best standard for hosting data online is through a
\textbf{Federated Data Repository}:

\begin{quote}
An online network of connected repositories that use similar standards
to collectively store data for discovery and access. Uploading a dataset
to one node of a repository will make it available through all other
nodes.
\end{quote}

Such repositories follow strict but necessary guidelines to ensure your
data live forever so long as the internet exists. The data are
definitely findable (e.g., through a web search), accessible (free to
download), and interoperable (accepted standards are ensured). The
``reproducible'' aspect can be debatable, but that can be solved through
other means (e.g., code sharing).

Some examples of data repositories, most are domain-specific:

\begin{itemize}
\tightlist
\item
  \href{https://knb.ecoinformatics.org/}{KNB}: Knowledge Network for
  Biocomplexity, a general purpose repository for ecological data
\item
  \href{https://www.hydroshare.org/}{HydroShare}: Data and models used
  in hydrology
\item
  \href{https://opc.dataone.org/data}{OPC}: California Ocean Protection
  Council, marine and coastal datasets
\end{itemize}

\hypertarget{code-of-conduct}{%
\subsection{Code of Conduct}\label{code-of-conduct}}

Every great collaborative team does not begin work before a Code of
Conduct is created. This documents a set of community and social
standards within which the work can be completed. It ensures all
viewpoints are heard and respected and establishes a means by which
conflicts can be resolved.

Here's a great example from our friends at
\href{https://www.openscapes.org/code-of-conduct/}{openscapes} and one
from the \href{https://ropensci.org/code-of-conduct/}{ROpenSci}
community. Some guidance for creating a code of conduct is also provided
in this
\href{https://ropensci.org/blog/2016/12/21/commcallv12-review-coc/}{blog
post} from ROpenSci.

The goal of every code of conduct is to ensure an agreed upon set of
norms are used by all team members to help create a safe and positive
experience.

\includegraphics[width=1em,height=1em]{./collaborate_files/figure-pdf/fa-icon-20d474448f872ee43905e611a2502347.pdf}
Exercise and discussion

Develop a code of conduct for your group in a shared workspace. Items to
consider:

\begin{enumerate}
\def\labelenumi{\arabic{enumi}.}
\tightlist
\item
  What are reasonable expectations for the group members?
\item
  How is inclusion defined and encouraged?
\item
  How are similarities and differences recognized?
\item
  How will conflicts be managed?
\end{enumerate}

\hypertarget{tools-for-collaboration}{%
\section{Tools for collaboration}\label{tools-for-collaboration}}

Now we introduce some specific web-based tools that you can use to
improve collaboration and openness. We present them as a suite of
options to consider based on the pros and cons associated with each
tool. This is by no means a comprehensive list, but it should get you
started towards better collaboration to leverage open science.

\hypertarget{slack}{%
\subsection{Slack}\label{slack}}

\includegraphics[width=0.15\textwidth,height=\textheight]{./img/slack-logo-icon.png}

\url{https://slack.com/}

\includegraphics[width=0.75em,height=1em]{./collaborate_files/figure-pdf/fa-icon-8447cf12acba1bc263badfefa79f6dd7.pdf}
What

An online messaging platform for internal communication. Conversations
can be organized by topic (via channels) or you can send direct messages
to one or more team members. You can have multiple workspaces for
different groups.

\includegraphics[width=1em,height=1em]{./collaborate_files/figure-pdf/fa-icon-e4ee65476be467d7be8a1ae9cb02ffda.pdf}
Pros

Alleviate email overload through quick, informal messaging. Offers a
fresh approach to online communication.

\includegraphics[width=1em,height=1em]{./collaborate_files/figure-pdf/fa-icon-c9ddf618df48250f2c0ed0cb59d4f20d.pdf}
Cons

Yet another thing to monitor. Free subscription limits archive of
messages. Communication is limited to only those individuals in a
workspace.

\hypertarget{trello}{%
\subsection{Trello}\label{trello}}

\includegraphics[width=0.15\textwidth,height=\textheight]{./img/trello.png}

\url{https://trello.com/}

\includegraphics[width=0.75em,height=1em]{./collaborate_files/figure-pdf/fa-icon-8447cf12acba1bc263badfefa79f6dd7.pdf}
What

A Kanban style workflow organization platform. Can be used for personal
organization or in teams. Card management allows you to assign due
dates, add attachments, make checklists, assign tasks to yourself or
team members, and label by themes.

\includegraphics[width=1em,height=1em]{./collaborate_files/figure-pdf/fa-icon-e4ee65476be467d7be8a1ae9cb02ffda.pdf}
Pros

Easy to use and can upgrade with ``power-ups'' for integration with
other services (e.g., Google). Use across locations (e.g., from home or
in the office) is easy because it's based in a web browser.

\includegraphics[width=1em,height=1em]{./collaborate_files/figure-pdf/fa-icon-c9ddf618df48250f2c0ed0cb59d4f20d.pdf}
Cons

Not entirely open because it's only visible to yourself or those you
explicitly invite. Free version is limited to only a handful of
``power-ups''.

\hypertarget{google-drive}{%
\subsection{Google Drive}\label{google-drive}}

\includegraphics[width=0.15\textwidth,height=\textheight]{./img/Google_Drive_logo.png}

\url{https://google.com/drive}

\includegraphics[width=0.75em,height=1em]{./collaborate_files/figure-pdf/fa-icon-8447cf12acba1bc263badfefa79f6dd7.pdf}
What

Cloud-based platform for sharing documents, worksheets, slides, etc.
Follows a familiar file-based GUI structure that is common to most
operating systems.

\includegraphics[width=1em,height=1em]{./collaborate_files/figure-pdf/fa-icon-e4ee65476be467d7be8a1ae9cb02ffda.pdf}
Pros

Easy to use and can be a very open space for collaboration. Fairly
interoperable with different file formats. Some functionality with
version control (i.e., ability to ``revert'' to previous versions and to
view changes).

\includegraphics[width=1em,height=1em]{./collaborate_files/figure-pdf/fa-icon-c9ddf618df48250f2c0ed0cb59d4f20d.pdf}
Cons

Requires a Google account and access can be tricky depending on
institution. Even though some versioning is provided, the format can
encourage poor file management. Who knows what Google is doing with your
data.

\hypertarget{office-365}{%
\subsection{Office 365}\label{office-365}}

\includegraphics[width=0.25\textwidth,height=\textheight]{./img/office365.png}

\url{https://www.microsoft.com/en-us/microsoft-365}

\includegraphics[width=0.75em,height=1em]{./collaborate_files/figure-pdf/fa-icon-8447cf12acba1bc263badfefa79f6dd7.pdf}
What

Cloud-based platform for secure sharing of Microsoft documents,
worksheets, slides, etc.

\includegraphics[width=1em,height=1em]{./collaborate_files/figure-pdf/fa-icon-e4ee65476be467d7be8a1ae9cb02ffda.pdf}
Pros

Easy to use and fully supports Microsoft products. Low barrier of
inclusion to others that are already using Microsoft products (should be
most folks).

\includegraphics[width=1em,height=1em]{./collaborate_files/figure-pdf/fa-icon-c9ddf618df48250f2c0ed0cb59d4f20d.pdf}
Cons

Requires a Microsoft account and access can be tricky depending on
institution. Maintains dependency on expensive Microsoft products that
don't facilitate reproducibility. Very often used in closed workflows.

\hypertarget{github}{%
\subsection{GitHub}\label{github}}

\includegraphics[width=0.25\textwidth,height=\textheight]{./img/Octocat.png}

\url{https://github.com}

\includegraphics[width=0.75em,height=1em]{./collaborate_files/figure-pdf/fa-icon-8447cf12acba1bc263badfefa79f6dd7.pdf}
What

Cloud-based platform for sharing code with Git version control. Supports
sharing of most file types, although code and text-based files are the
primary use.

\includegraphics[width=1em,height=1em]{./collaborate_files/figure-pdf/fa-icon-e4ee65476be467d7be8a1ae9cb02ffda.pdf}
Pros

Collaborative and fully transparent work environment for files under
version control. Supports workflow management through issue tracking and
Kanban style project boards. Links to third-party platforms for
archiving and DOI generation (e.g., \href{https://zenodo.org/}{Zenodo}).
Octocat mascot is super cute.

\includegraphics[width=1em,height=1em]{./collaborate_files/figure-pdf/fa-icon-c9ddf618df48250f2c0ed0cb59d4f20d.pdf}
Cons

Learning curve is steep if you want to fully leverage version control.
Not a formal data archival service by itself and file sizes are limited.

\includegraphics[width=1.25em,height=1em]{./collaborate_files/figure-pdf/fa-icon-5a0e5cd7bfbf585f8546c1ff4cebae73.pdf}
Watch and learn

Setting up a project management system with GitHub.

\includegraphics[width=1em,height=1em]{./collaborate_files/figure-pdf/fa-icon-20d474448f872ee43905e611a2502347.pdf}
Exercise and discussion

In small groups, setup a shared workspace using GitHub and create a
project management board.

\hypertarget{balancing-and-managing-expectations}{%
\section{Balancing and managing
expectations}\label{balancing-and-managing-expectations}}

Adopting new workflows and tools to facilitate collaboration are steps
in the right direction for open science. Clearly, change from
established norms and practices is not easy. It's important to develop a
realistic expectation of how this change may play out in the wild.

First, realize that these changes are not going to happen overnight.
Breaking free of entrenched workflows is like quitting bad habits. It's
not easy because they're comfortable, familiar, and often habitual.
Start small and work gradually towards adopting new ideas. Incremental
progress is the name of the game.

Getting others on board is another serious challenge. The sources of
frustration you might have at the personal level apply to anyone else
working in your team. Even more so, institutional roadblocks may exist.
New software can raise red flags for IT support staff. It's important to
work with them to develop trust for the software and an accepted process
for installation.

Finally, the tools above are temporary. The single constant in open
science is change. Use the tools with full realization that they may be
appropriate for now, but something (hopefully) better will replace it in
the future.

\bookmarksetup{startatroot}

\hypertarget{impact}{%
\chapter{Open science for impactful products}\label{impact}}

\hypertarget{lesson-outline-2}{%
\section{Lesson Outline}\label{lesson-outline-2}}

\begin{itemize}
\tightlist
\item
  \protect\hyperlink{goals-and-motivation-3}{Goals and motivation}
\item
  \protect\hyperlink{data-as-the-foundation-for-open-science}{Data as
  the foundation for open science}
\item
  \protect\hyperlink{tidy}{Principles of tidy data}
\item
  \protect\hyperlink{data-dictionaries}{Data dictionaries}
\item
  \protect\hyperlink{importance-of-metadata}{Importance of metadata}
\end{itemize}

\hypertarget{goals-and-motivation-2}{%
\section{Goals and motivation}\label{goals-and-motivation-2}}

This is the third module in our four hour workshop on open science. Now
we focus on core principles for data management as the foundation for
open science. We discuss the role of data management to support
decisions using open science. Then, we introduce the concepts of tidy
data as a unified format for storing information. We close with a
discussion of metadata concepts and tools to make sure your data have a
life beyond the project.

\begin{itemize}
\tightlist
\item
  \textbf{Goal}: understand best practices for data management as a key
  concept for open science
\item
  \textbf{Motivation}: cultivate data as a living, shared resource
\end{itemize}

\hypertarget{data-as-the-foundation-for-open-science}{%
\section{Data as the foundation for open
science}\label{data-as-the-foundation-for-open-science}}

In the last module, we talked about collaboration as the single most
important activity of open science. So, why are we now talking about
data management? Understanding the tools of collaboration allows you to
better engage with your colleagues and partners, but open engagement
will mean nothing if your data look like garbage.

\begin{figure}

{\centering \includegraphics[width=4in,height=\textheight]{./img/datameme.jpg}

}

\end{figure}

You can probably recall several past instances when poor data management
has been a challenge for open collaboration. Here are a few real-world
examples:

\begin{enumerate}
\def\labelenumi{\arabic{enumi}.}
\tightlist
\item
  A collaborator calls you on the phone asking about a historical
  dataset from an old report. You spend several hours tracking down this
  information because you don't know where it is. The data you
  eventually find and provide to your collaborator has no documentation
  and they don't know how to use it or use it inappropriately.
\item
  You receive a deliverable from a project partner that was stipulated
  in a scope of work. This deliverable comes in multiple formats with no
  reproducible workflow to recreate the datasets. You are unable to
  verify the information, eroding your faith in the final product and
  making it impossible to update the results in the future.
\item
  An annual reporting product requires using new data each year. The
  staff member in charge of this report spends several days gathering
  the new data and combining it with the historical data. Other projects
  are on hold until this report is updated. Stakeholders that use this
  report to make decisions do not trust or misunderstand the product
  because the steps for its creation are opaque.
\end{enumerate}

Anybody who has ever worked with data knows that it comes in many shapes
and sizes, most more like the right side of the above picture. Poor data
management occurs for many reasons, but here's a few of the common
reasons:

\begin{itemize}
\tightlist
\item
  What's easy for recording data doesn't usually translate to easy
  analysis
\item
  Egregious use of Excel as data management software
\item
  Metadata is a chore that is often an afterthought
\end{itemize}

It's often said that 90\% of working with data is cleaning (or
``wrangling''), whereas the actual analysis and interpretation part is a
small fraction of your total effort. Using better data management
practices will not only help you save time, it's also a service for your
colleagues and future collaborators. Therefore, better data management
leads to more open science.

The \protect\hyperlink{fair}{FAIR} principles outlined in the first
module are especially useful when working with data for open science
applications. Many of the collaborative tools in the second module can
help you work towards putting FAIR data into practice. In this module,
we'll go a step further to discuss how data structure, including
metadata, can produce a FAIR dataset.

\hypertarget{tidy}{%
\section{Principles of tidy data}\label{tidy}}

At their core, tabular data allow you to store information as
observations in rows and variables in columns, yet it's very common to
try to make a data table more than it should be. Unless you spend a lot
of time working with data, it can be difficult to recognize common
mistakes that lead to \emph{table abuse}.

Before we get into tidy data, we need to discuss some of the downfalls
of Excel as a data management system. There are
\href{http://www.eusprig.org/horror-stories.htm}{many examples} that
demonstrate how Excel has contributed to costly mistakes through the
abuse of tables, often to the detriment of science (Ziemann, Eren, and
El-Osta 2016).

\begin{figure}

{\centering \includegraphics[width=0.7\textwidth,height=\textheight]{./img/excel_dates.JPG}

}

\end{figure}

Excel allows you to abuse your data in many ways, such as adding color
to cells, embedding formulas, and automatically formatting cell types.
The problem occurs when this organization becomes ambiguous and only has
meaning inside the head of the person who created the spreadsheet.
Embedding formulas that reference specific locations in or across
spreadsheets is also a nightmare scenario for reproducibility.

\begin{figure}

{\centering \includegraphics[width=0.8\textwidth,height=\textheight]{./img/excel_bad.png}

}

\end{figure}

If you absolutely must use Excel to store data, the only acceptable
format is a rectangular, flat file. What do we mean by this?

A \textbf{rectangular} file:

\begin{quote}
Store data only in rows and columns in matrix format (e.g., 10 rows x 5
columns), with no ``dangling'' cells that have values outside of the
grid or more than one table in a spreadsheet.
\end{quote}

A \textbf{flat} file:

\begin{quote}
No cell formatting, no embedded formulas, no multiple spreadsheets in
the same file, and data entered only as alphanumeric characters.
\end{quote}

Broman and Woo (2018) provide an excellent guide that expands on these
ideas. Essentially, these best practices force you to isolate the
analysis from the data - many people use Excel to mix the two, leading
to problems.

Now we can talk about tidy data. The ``tidy'' data principles developed
by Hadley Wickham (Wickham 2014) are a set of simple rules for storing
tabular data that have motivated the development of the wildly popular
\href{https://www.tidyverse.org/}{tidyverse} suite of R packages
(Wickham et al. 2019). The rules are simple:

\begin{enumerate}
\def\labelenumi{\arabic{enumi}.}
\tightlist
\item
  Each variable must have its own column;
\item
  Each observation must have its own row; and,
\item
  Each value must have its own cell.
\end{enumerate}

Graphically, these rules are shown below (from Wickham and Grolemund
2017):

\includegraphics[width=6.4in,height=\textheight]{./img/tidy-1.png}

Using these principles for data storage may seem unnatural at first
because of a difference between what's easy for entering data versus
what makes sense for downstream analyses. For example, dates are often
spread across multiple columns, such as having one column for each year
of data where the header indicates the year that applies to data in each
column.

Using a tidy format also allows you to more easily join data between
tables. This is a common task when you have information spread between
different tables because: 1) it might not make sense to keep the data in
the same table, and 2) the analysis depends on information from both
tables. Tidy data shared between tables can be linked using a ``key'' as
a common identifier.

\includegraphics[width=6.26in,height=\textheight]{./img/joins.png}

\includegraphics[width=1.25em,height=1em]{./impact_files/figure-pdf/fa-icon-5a0e5cd7bfbf585f8546c1ff4cebae73.pdf}
Watch and learn

Making an
\href{https://github.com/tbep-tech/cerf-os-workshop/raw/master/data/untidy.xlsx}{untidy}
dataset
\href{https://github.com/tbep-tech/cerf-os-workshop/raw/master/data/tidy.xlsx}{tidy}
using Excel.

\includegraphics[width=1.25em,height=1em]{./impact_files/figure-pdf/fa-icon-5a0e5cd7bfbf585f8546c1ff4cebae73.pdf}
Watch and learn

Making an
\href{https://github.com/tbep-tech/cerf-os-workshop/raw/master/data/untidy.xlsx}{untidy}
dataset
\href{https://github.com/tbep-tech/cerf-os-workshop/raw/master/data/tidy.xlsx}{tidy}
using R.

\begin{Shaded}
\begin{Highlighting}[]
\FunctionTok{library}\NormalTok{(readxl)}
\FunctionTok{library}\NormalTok{(dplyr)}
\FunctionTok{library}\NormalTok{(tidyr)}

\CommentTok{\# import and wrangle}
\NormalTok{dat }\OtherTok{\textless{}{-}} \FunctionTok{read\_excel}\NormalTok{(}\StringTok{\textquotesingle{}data/untidy.xlsx\textquotesingle{}}\NormalTok{, }\AttributeTok{skip =} \DecValTok{1}\NormalTok{) }\SpecialCharTok{\%\textgreater{}\%} 
  \FunctionTok{fill}\NormalTok{(Location) }\SpecialCharTok{\%\textgreater{}\%} 
  \FunctionTok{pivot\_longer}\NormalTok{(}\AttributeTok{cols =} \StringTok{\textasciigrave{}}\AttributeTok{2019}\StringTok{\textasciigrave{}}\SpecialCharTok{:}\StringTok{\textasciigrave{}}\AttributeTok{2021}\StringTok{\textasciigrave{}}\NormalTok{, }\AttributeTok{names\_to =} \StringTok{\textquotesingle{}Year\textquotesingle{}}\NormalTok{, }\AttributeTok{values\_to =} \StringTok{\textquotesingle{}acres/category\textquotesingle{}}\NormalTok{) }\SpecialCharTok{\%\textgreater{}\%} 
  \FunctionTok{separate}\NormalTok{(}\AttributeTok{col =} \StringTok{\textasciigrave{}}\AttributeTok{acres/category}\StringTok{\textasciigrave{}}\NormalTok{, }\AttributeTok{into =} \FunctionTok{c}\NormalTok{(}\StringTok{\textquotesingle{}acres\textquotesingle{}}\NormalTok{, }\StringTok{\textquotesingle{}category\textquotesingle{}}\NormalTok{), }\AttributeTok{sep =} \StringTok{\textquotesingle{}/\textquotesingle{}}\NormalTok{)}
\NormalTok{dat}
\end{Highlighting}
\end{Shaded}

\includegraphics[width=1em,height=1em]{./impact_files/figure-pdf/fa-icon-20d474448f872ee43905e611a2502347.pdf}
Exercise and discussion

Now try it on your own. Download the
\href{https://github.com/tbep-tech/cerf-os-workshop/raw/master/data/untidy.xlsx}{untidy}
dataset and make it tidy using your preferred software.

\hypertarget{data-dictionaries}{%
\section{Data dictionaries}\label{data-dictionaries}}

Once you understand the tidy principles, you'll find that analysis is
much, much easier. More importantly, this also facilitates documentation
for the explicit purpose of more open sharing. As responsible data
stewards, we need to think of data as a living resource for open
collaboration that is not just a means for more conventional scientific
products (e.g., a publication). Data are increasingly being documented
and cited as unique entities and we need proper documentation to help
support the FAIR principles.

A good first step in documentation is to create a data dictionary. This
will help us when we start creating
\protect\hyperlink{importance-of-metadata}{metadata}. Think of this as
the specific description of the contents of a tabular data file.
Developing a data dictionary not only helps with metadata, but also
helps you think more clearly about your data.

A data dictionary describes column names and the type of data in each
column. Simple things like how you name a data column can have larger
implications for downstream analysis pipelines or interpretability of a
dataset.

Here's an example of a data dictionary for a made up dataset. Without
seeing the actua data, you get a good sense of what's included and
acceptable values for adding new data.

\begin{figure}

{\centering \includegraphics[width=0.8\textwidth,height=\textheight]{./img/datadict.png}

}

\end{figure}

Here we provide some general guidelines for developing your own data
dictionary.

\begin{itemize}
\item
  Column names

  \begin{itemize}
  \tightlist
  \item
    Be as descriptive as possible while trying to keep the name as short
    as possible. Really long names with lots of detail can be just as
    frustrating as very short names with very little detail. The column
    name should be intuitive to point the analyst in the right
    direction.\\
  \item
    Try to avoid spaces or commas in column names since some software
    may interpret that as the start of a new column.\\
  \item
    It may also be useful to identify a ``plot name'' for each column
    name that uses proper spelling, punctuation, and units (if
    applicable). Using a column name directly in a graphic is generally
    a bad idea since their meaning outside of the dataset may not be
    obvious.
  \end{itemize}
\item
  Data types

  \begin{itemize}
  \tightlist
  \item
    Describe the type of data in each column, e.g., numerical
    measurements, categorical descriptors, or counts of observations.
    Never, ever mix data types in the same column.\\
  \item
    If your data are continuous numeric values, try to identify an
    acceptable range for the values, e.g., are there minimum or maximum
    values that would indicate the data are out of range? Also make note
    of the units that were used.\\
  \item
    For categorical descriptors, identify all possible categories that
    are acceptable values for the column, e.g., small, medium, or large
    for a qualitative descriptor of size.\\
  \item
    For dates, make note of the format, e.g., YYYY-MM-DD. For time,
    identify the timezone.
  \end{itemize}
\end{itemize}

\includegraphics[width=1em,height=1em]{./impact_files/figure-pdf/fa-icon-20d474448f872ee43905e611a2502347.pdf}
Exercise and discussion

Create a data dictionary for the
\href{https://github.com/tbep-tech/cerf-os-workshop/raw/master/data/tidy.xlsx}{tidy}
dataset from the previous example.

\hypertarget{importance-of-metadata}{%
\section{Importance of metadata}\label{importance-of-metadata}}

How many times have you been sent a dataset without any idea what it
contains or why it was created? How are you sure the information is
valid and that your analysis takes into account the limitations of the
data? How many times have you willfully sent someone a dataset without
fully providing this information?

Without metadata it's impossible to know critical details about a
dataset that can inform its analysis, and more importantly, its use to
inform decision-making. Curating data should be synonymous with metadata
generation and is an important part of open science. We cannot provide
open data in good faith without also providing metadata.

Metadata is literally defined as ``data about data''. It varies from
simple text descriptions of a dataset, such as ``who'', ``what'',
``when'', ``where'', ``why'', and ``how'', to more formalized standards
with the intent of preparing your data for archival in a long-term
repository. A useful definition is provided by Gilliland (2016):

\begin{quote}
A suite of industry or disciplinary standards as well as additional
internal and external documentation and other data necessary for the
identification, representation, interoperability, technical management,
performance, and use of data contained in an information system.
\end{quote}

Why don't we see more metadata in the wild? Short answer is that it's
often an afterthought, if considered at all. Creating metadata is
usually tedious and the return on investment is not apparent at onset of
a project. However, the collective growth of sciences and its
application to real world problems is dependent on metadata.

The US Geological Survey provides a useful document on creating
\href{https://prd-wret.s3.us-west-2.amazonaws.com/assets/palladium/production/atoms/files/Metadata\%20in\%20Plain\%20Language_508compliant.pdf}{Metadata
in ``plain language''} to distill the basic information contained in a
metadata file. It provides a workflow for answering the ``who'',
``what'', ``when'', ``where'', ``why'', and ``how'' questions for
metadata. Below is a brief synopsis:

\begin{quote}
What does the dataset describe?
\end{quote}

Information here would include very basic details about the dataset
including a \textbf{title}, \textbf{geographic extent}, and
\textbf{period of time} covered by the data. For geographic extent, this
may often include explicit coordinates covering the study area. Location
is useful for indexing your dataset relative to others, if for example,
a researcher wanted to find data for all studies in the geographic
extent of Tampa Bay.

\begin{quote}
Who produced the dataset?
\end{quote}

This would be yourself and anyone else who has made a significant
contribution to the development of a dataset. Data are increasingly
being used as citable resources and including individuals that were
important in its generation ensures proper attribution. If someone has
spent hours toiling in the field to collect the data or hours visually
scanning a spreadsheet for quality control, include them!

\begin{quote}
Why was the dataset created?
\end{quote}

Describing why a dataset was created is critically important for
understanding context. If others want to use your data, they need to
know if it's appropriate for their needs. Here you would describe the
goal or objectives of the research for which the data were collected. It
should be clear if there are limitations in your data defined by your
goals.

\begin{quote}
How was the dataset created?
\end{quote}

Here you would describe the methods used to generate the data, e.g.,
field sampling techniques, laboratory methods, etc. This information is
important so others can know if you've used proper and accepted methods
for generating the data. Citing existing SOPs or methods that are
recognized standards in your field would be appropriate.

\begin{quote}
How reliable are the data?
\end{quote}

It's also important to explicitly note instances when the data could be
questionable or inappropriate to use. Here you could describe any
quality assurance or quality control (QAQC) checks that were used on the
data. There are often formalized ways to do so, such as codes or
descriptors in tabular data defining QAQC values (e.g., data in range,
below detection limits, sensor out of service, etc.).

\begin{quote}
How can someone get a copy of the dataset?
\end{quote}

Good metadata has information on who contact for getting the data. This
contact may not be the same as who created the dataset (e.g., IT staff).
For archived or publicly available data, this information is more
important for who to contact should someone have questions. Information
on obtaining a copy of the data should also describe any special
software or licensing issues related to accessing the data.

Once you've gathered this information, how do you turn it into literal
metadata? It depends on how deep you want to go. At it's simplest, your
metadata could be a simple text file with answers to the questions. Or
it could be a specific file format used by modern data archive
repositories (e.g., \href{https://docs.ropensci.org/EML/}{EML} format).

Here's an example of a bare bones metadata file. One could easily type
this up in a text file or spreadsheet.

\includegraphics[width=6.17in,height=\textheight]{./img/desotometa.PNG}

\includegraphics[width=1em,height=1em]{./impact_files/figure-pdf/fa-icon-20d474448f872ee43905e611a2502347.pdf}
Exercise and discussion

Create a metadata file for the
\href{https://github.com/tbep-tech/cerf-os-workshop/raw/master/data/tidy.xlsx}{tidy}
dataset from the previous example. Just start by using your data
dictionary and develop narrative answers for ``who'', ``what'',
``when'', ``where'', ``why'', and ``how'' to describe your dataset. Get
creative with your descriptions since this is a made up dataset.

\bookmarksetup{startatroot}

\hypertarget{implement}{%
\chapter{Lowering barriers to inclusion and addressing key
critiques}\label{implement}}

\hypertarget{lesson-outline-3}{%
\section{Lesson Outline}\label{lesson-outline-3}}

\begin{itemize}
\tightlist
\item
  \protect\hyperlink{goals-and-motivation-3}{Goals and motivation}
\item
  \protect\hyperlink{learning-curves}{Learning curves}
\item
  \protect\hyperlink{fear-of-exposure}{Fear of exposure}
\item
  \protect\hyperlink{what-does-it-mean-to-be-open}{What does it mean to
  be open?}
\item
  \protect\hyperlink{something-is-better-than-nothing}{Something is
  better than nothing}
\end{itemize}

\hypertarget{goals-and-motivation-3}{%
\section{Goals and motivation}\label{goals-and-motivation-3}}

We finish our workshop with a discussion of what it means to use open
science in the real world. It's great to talk about the value of open
science and the tools you can use, but it's a completely different ball
game when it comes to putting these ideas into practice. Our goal is
that you leave this workshop an advocate and early adopter for the ideas
we discussed today - spread these ideas to your peers and colleagues! To
realistically achieve this goal, we will talk about some of the
challenges you will face so you can develop a realistic expectation of
what's to come.

\begin{itemize}
\tightlist
\item
  \textbf{Goal}: Understand common hurdles in adopting open science and
  how to overcome them\\
\item
  \textbf{Motivation}: Become the ``open science'' expert at your
  institution!
\end{itemize}

\hypertarget{learning-curves}{%
\section{Learning curves}\label{learning-curves}}

\begin{figure}

{\centering \includegraphics[width=0.6\textwidth,height=\textheight]{./img/challenge.png}

}

\end{figure}

\includegraphics[width=1em,height=1em]{./implement_files/figure-pdf/fa-icon-e38cc40f3c6469ef5f44a7764a5c5910.pdf}
Challenge

\emph{It's hard to learn new tools!}

\includegraphics[width=1em,height=1em]{./implement_files/figure-pdf/fa-icon-e4ee65476be467d7be8a1ae9cb02ffda.pdf}
Solution

\emph{It's an investment, look to the community!}

You've probably seen a graphic like this if you've ever taken a course
in R or Python. The hope is that you're able to quickly reach the land
of sunshine and bunnies, but the path is treacherous and even
insurmountable for some.

A huge obstacle in using open science is that the toolsets can have
steep learning curves. More popular platforms, such as Excel, are used
by many because they're simple and intuitive. However,
\protect\hyperlink{tidy}{as noted earlier}, FAIR workflows and tools are
sacrificed for ease of use.

Although it's true that adopting new tools will slow forward progress,
this is only temporary. Consider your path towards learning new
platforms an investment in your future. The immediate benefit may not be
apparent, but you'll soon wonder how you ever got by before.

It's also helpful to think about the broader community that can support
you along this journey. Learning alone can be discouraging and we
strongly recommend that you tap into the diverse community of educators,
mentors, bloggers, and friends that can help. Even you can create a
community of practice!

\begin{figure}

{\centering \includegraphics[width=0.75\textwidth,height=\textheight]{./img/code_hero.jpg}

}

\end{figure}

\includegraphics[width=1em,height=1em]{./implement_files/figure-pdf/fa-icon-20d474448f872ee43905e611a2502347.pdf}
Exercise and discussion

How can you engage your peers to develop a shared workspace to learn new
tools? What tools will you learn?

\hypertarget{fear-of-exposure}{%
\section{Fear of exposure}\label{fear-of-exposure}}

\begin{figure}

{\centering \includegraphics[width=0.3\textwidth,height=\textheight]{./img/fishbowl.png}

}

\end{figure}

\includegraphics[width=1em,height=1em]{./implement_files/figure-pdf/fa-icon-e38cc40f3c6469ef5f44a7764a5c5910.pdf}
Challenge

\emph{Being open makes me nervous!}

\includegraphics[width=1em,height=1em]{./implement_files/figure-pdf/fa-icon-e4ee65476be467d7be8a1ae9cb02ffda.pdf}
Solution

\emph{Being open helps you collaborate, increases competitiveness, and
creates a better scientific product!}

Practicing open science can feel like science in a fish bowl. Although
this is kind of the goal, many view this transparency as a liability.
Many fear having their ideas ``scooped'' or losing credibility because
of greater exposure of mistakes. These are real concerns that require
consideration when working towards more open workflows.

In conventional academic settings, competition for resources (e.g., via
grant funding) is a real issue and being open can be seen as a risk to
the competitive edge. We cannot dismiss this fact, but rather we can
think about a lack of openness as a hindrance to forward progress and
stifled creativity.

Think about being open as a means to finding your next collaborator.
Creating \protect\hyperlink{fair}{FAIR} data opens the door for others
to engage with your science. In fact, being open can increase the
competitiveness of research proposals by building a stronger team that
collaborates and shares data through better workflows.

First time practitioners of open science also worry about the risk of
``airing their dirty laundry''. By exposing the process and potential
mistakes, many worry that their integrity as scientists may be
questioned.

These fears are unfounded as the scientific process by definition is
iterative. Hypotheses are supported or refuted through trial and error -
if you're getting your answer after one pass, you're probably not doing
it right. Making the process more transparent can help build trust as
your collaborators can better appreciate how decisions and conclusions
were made.

Mistakes in research are also very common, much more so than many people
realize. By being open, it is true that mistakes are more visible, but
this also provides a mechanism for fixing. Being open can lead to a
better product by simply having more eyes on the process. It also helps
normalize mistakes as part of the process - perfection is an unrealistic
expectation.

\includegraphics[width=1em,height=1em]{./implement_files/figure-pdf/fa-icon-20d474448f872ee43905e611a2502347.pdf}
Exercise and discussion

What are your personal concerns about adopting open science?

\hypertarget{what-does-it-mean-to-be-open}{%
\section{What does it mean to be
open?}\label{what-does-it-mean-to-be-open}}

\includegraphics[width=1em,height=1em]{./implement_files/figure-pdf/fa-icon-e38cc40f3c6469ef5f44a7764a5c5910.pdf}
Challenge

\emph{People and institutions define open differently!}

\includegraphics[width=1em,height=1em]{./implement_files/figure-pdf/fa-icon-e4ee65476be467d7be8a1ae9cb02ffda.pdf}
Solution

\emph{Understand the context and demonstrate the value!}

Also realize that open science can mean different things to different
people. By extension, this also applies to institutions. We presented
the \protect\hyperlink{schools}{five schools} of open science to help
conceptualize ideas and tools when we discuss what it means to different
groups.

Think about your employer and what they might care about if you advocate
for adoption of open science. Do you need to convince them that there is
value in being open? What is their value proposition? What are the
hurdles to achieving openness at your institution?

For many institutions, being open may come with IT hurdles as you push
for alternative software platforms. Working with IT staff to develop
trust and comfort for new software may be your burden, but as always,
it's an investment in the future.

Maybe there are legal contexts to being open. For example, Florida has
\href{https://myfloridalegal.com/pages.nsf/Main/DC0B20B7DC22B7418525791B006A54E4}{the
``Sunshine'' law} that makes all government communications public
record. What does this mean for using new workflows in open science? Is
this is an improvement or a liability (see
\protect\hyperlink{fear-of-exposure}{previous section})?

If you're an administrator or manager, maybe you're the one that makes
the call about being open. It's important for you to create a culture
that promotes and supports open science. Allow space and time for your
staff to learn new skills. Realize that investing time in open science
is an investment in the future.

\includegraphics[width=1em,height=1em]{./implement_files/figure-pdf/fa-icon-20d474448f872ee43905e611a2502347.pdf}
Exercise and discussion

What does being open mean to you? What do you think being open means to
your employer?

\hypertarget{something-is-better-than-nothing}{%
\section{Something is better than
nothing}\label{something-is-better-than-nothing}}

\includegraphics[width=1em,height=1em]{./implement_files/figure-pdf/fa-icon-e38cc40f3c6469ef5f44a7764a5c5910.pdf}
Challenge

\emph{Doing all the things is impossible!}

\includegraphics[width=1em,height=1em]{./implement_files/figure-pdf/fa-icon-e4ee65476be467d7be8a1ae9cb02ffda.pdf}
Solution

\emph{Start small, incremental progress is the name of the game!}

First time open science enthusiasts can be overwhelmed by the apparent
need to check all the boxes on the open science list. There's often a
prevailing sentiment that you're not doing open science unless you do
all the things. This is simply not true. Just remember that doing
something is a huge improvement over doing nothing.

Openness in science exists on a spectrum. Your goal should be
incremental movement away from the completely closed end of the
spectrum. Perhaps you set a goal of only accomplishing one open science
task for a particular project. Maybe you start by developing a simple
metadata text file or developing a data dictionary. Or maybe you make a
commitment to try a new communication platform for collaborative
engagement.

Channeling this concept, Wilson et al. (2017) discuss ``good enough
practices'' in scientific computing, acknowledging that very few of us
are professionally trained in these disciplines and sometimes ``good
enough'' is all we can ask for. Lowenberg et al. (2021) also advocate
for simple adoption, rather than perfection, when it comes to data
citation practices.

So, be kind to yourself when learning new skills and realize that the
first step will likely be frustration, but through frustration comes
experience. The more comfortable you become with a new task, the more
likely you'll be able to attempt additional tasks in the future.

\includegraphics[width=1em,height=1em]{./implement_files/figure-pdf/fa-icon-20d474448f872ee43905e611a2502347.pdf}
Exercise and discussion

What are some simple things you can do to begin adopting open science?

\bookmarksetup{startatroot}

\hypertarget{resources-for-continued-learning}{%
\chapter*{Resources for continued
learning}\label{resources-for-continued-learning}}
\addcontentsline{toc}{chapter}{Resources for continued learning}

The following is a non-exhaustive list of additional resources you can
use for continued learning on your open science journey.

\hypertarget{open-science-websites}{%
\subsection*{Open Science Websites}\label{open-science-websites}}
\addcontentsline{toc}{subsection}{Open Science Websites}

\begin{itemize}
\tightlist
\item
  \href{https://nceas.github.io/oss-2017/lessons.html}{NCEAS Open
  Science for Synthesis 2017 workshop}
\item
  \href{https://learning.nceas.ucsb.edu/2020-02-RRCourse/}{NCEAS
  Reproducible Research Techniques}
\item
  \href{https://osf.io/k75r9/}{Open Science Foundation open science
  workshop}
\item
  \href{https://www.openscapes.org/}{Openscapes}
\item
  \href{https://openscapes.github.io/series/}{Openscapes Champions
  Lesson Series}
\item
  \href{https://www.nature.com/articles/d41586-019-03335-4}{Supercharge
  your research: A 10 week plan for open data science}
\item
  \href{https://ropensci.org/blog/2016/12/21/commcallv12-review-coc/}{ROpenSci
  guidance on creating a Code of Conduct}
\item
  \href{https://noaa-iea.github.io/r3-train/index.html}{NOAA
  Reproducible Reporting with R}
\item
  \href{https://peerj.com/collections/50-practicaldatascistats}{PeerJ
  collection on practical data science}
\end{itemize}

\hypertarget{data-management-tools}{%
\subsection*{Data Management Tools}\label{data-management-tools}}
\addcontentsline{toc}{subsection}{Data Management Tools}

\begin{itemize}
\tightlist
\item
  \href{https://environmentaldatainitiative.org/dm-resources/}{Environmental
  Data Initiative Data Management Resources}
\item
  \href{https://dmptool.org/}{University of California DMPTool}
\item
  \href{https://www.usgs.gov/products/data-and-tools/data-management/metadata-creation}{US
  Geological Survey resources for Metadata Creation}
\item
  \href{https://ds-wizard.org/}{ELIXIR and others Data Stewardship
  Wizard}
\item
  \href{https://tbep-tech.github.io/data-management-sop}{TBEP Data
  Management SOP}
\end{itemize}

\hypertarget{tbep-r-trainings}{%
\subsection*{TBEP R Trainings}\label{tbep-r-trainings}}
\addcontentsline{toc}{subsection}{TBEP R Trainings}

\begin{itemize}
\tightlist
\item
  \href{https://tbep-tech.github.io/pep-r-training}{Peconic Estuary
  Program R training},
  \href{https://drive.google.com/file/d/1ZjVHFrVpw2uTKZw-BmD29umdbl6viutM/view?usp=sharing}{recording}
\item
  \href{https://tbep-tech.github.io/ecometab-r-training}{NERRS ecosystem
  metabolism R training}
  \href{https://tbep-tech.github.io/tbep-r-training/}{TBEP June 2020 R
  training},
  \href{https://www.youtube.com/watch?v=_RI4XMRWeV0\&list=PLfJ6-D-exF9RM5TPtT4T0nxieqpr_R4pJ}{recordings}
\item
  \href{https://tbep-tech.github.io/tbep-r-training/functions.html}{Writing
  functions in R}
\item
  \href{https://tbep-tech.github.io/tbep-r-training/packages.html}{R
  package development workflow}
\item
  \href{https://tbep-tech.github.io/tbep-r-training/shiny.html}{A soft
  introduction to Shiny}
\end{itemize}

\hypertarget{r-lessons-tutorials}{%
\subsection*{R Lessons \& Tutorials}\label{r-lessons-tutorials}}
\addcontentsline{toc}{subsection}{R Lessons \& Tutorials}

\begin{itemize}
\tightlist
\item
  \href{http://data-lessons.github.io/gapminder-R/}{Software Carpentry:
  R for Reproducible Scientific Analysis}
\item
  \href{https://datacarpentry.org/geospatial-workshop/}{Data Carpentry:
  Geospatial Workshop}
\item
  \href{http://www.datacarpentry.org/R-ecology-lesson/}{Data Carpentry:
  R for Data Analysis and Visualization of Ecological Data}
\item
  \href{http://www.datacarpentry.org/spreadsheet-ecology-lesson/}{Data
  Carpentry: Data Organization in Spreadsheets}
\item
  \href{https://www.r4wrds.com/}{R for Water Resources Data Science}
\item
  \href{https://www.rstudio.com/resources/webinars/}{RStudio Webinars,
  many topics}
\item
  \href{http://rforcats.net/}{R For Cats: Basic introduction site, with
  cats!}
\item
  \href{http://www.rstudio.com/resources/cheatsheets/}{Topical
  cheatsheets from RStudio, also viewed from the help menu}
\item
  \href{http://cran.r-project.org/doc/contrib/Short-refcard.pdf}{Cheatsheet
  from CRAN of base R functions}
\item
  \href{https://github.com/allisonhorst/stats-illustrations}{Totally
  awesome R-related artwork by Allison Horst}
\item
  \href{http://www.stat.columbia.edu/~tzheng/files/Rcolor.pdf}{Color
  reference PDF with text names},
  \href{https://www.nceas.ucsb.edu/sites/default/files/2020-04/colorPaletteCheatsheet.pdf}{Color
  cheatsheet PDF from NCEAS}
\end{itemize}

\hypertarget{r-ebookscourses}{%
\subsection*{R eBooks/Courses}\label{r-ebookscourses}}
\addcontentsline{toc}{subsection}{R eBooks/Courses}

\begin{itemize}
\tightlist
\item
  \href{http://stat545.com/}{Jenny Bryan's Stat545.com}
\item
  \href{http://r4ds.had.co.nz/}{Garrett Grolemund and Hadley Wickham's R
  For Data Science}
\item
  \href{https://ismayc.github.io/moderndiver-book/}{Chester Ismay and
  Albert Y. Kim's Modern DiveR}
\item
  \href{http://tidytextmining.com/}{Julia Silge and David Robinson Text
  Mining with R}
\item
  \href{http://adv-r.had.co.nz/}{Hadley Wickham's Advanced R}
\item
  \href{https://r4ds.had.co.nz}{Hadley Wickham's R for Data Science}
\item
  \href{https://bookdown.org/yihui/rmarkdown/}{Yihui Xie R Markdown: The
  Definitive Guide}
\item
  \href{http://www.cookbook-r.com/}{Winston Chang R Graphics Cookbook}
\item
  \href{http://book.ecosens.org/RSEbook/}{Wegman et al.~Remote Sensing
  and GIS for Ecologists: Using Open Source Software}
\item
  \href{https://geocompr.robinlovelace.net/}{Lovelace et
  al.~Geocomputation with R}
\item
  \href{https://keen-swartz-3146c4.netlify.app/index.html}{Edszer
  Pebesma and Roger Bivand Spatial Data Science}
\end{itemize}

\hypertarget{gitgithub}{%
\subsection*{Git/Github}\label{gitgithub}}
\addcontentsline{toc}{subsection}{Git/Github}

\begin{itemize}
\tightlist
\item
  \href{http://happygitwithr.com/}{Jenny Bryan's Happy Git and Github
  for the useR}
\item
  \href{https://ourcodingclub.github.io/2017/02/27/git.html}{Coding Club
  Intro to Github}
\end{itemize}

\bookmarksetup{startatroot}

\hypertarget{references}{%
\chapter*{References}\label{references}}
\addcontentsline{toc}{chapter}{References}

\hypertarget{refs}{}
\begin{CSLReferences}{1}{0}
\leavevmode\vadjust pre{\hypertarget{ref-Beck20}{}}%
Beck, M. W., C. O'Hara nad J. S. S. Lowndes, R. D. Mazor, S. Theroux, D.
J. Gillett, B. Lane, and G. Gearheart. 2020. {``The Importance of Open
Science for Biological Assessment of Aquatic Environments.''}
\emph{PeerJ} 8: e9539.
https://doi.org/\url{https://doi.org/10.7717/peerj.9539}.

\leavevmode\vadjust pre{\hypertarget{ref-Broman18}{}}%
Broman, K. W., and K. H. Woo. 2018. {``Data Organization in
Spreadsheets.''} \emph{The American Statistician} 72 (1): 2--10.
\url{https://doi.org/10.1080/00031305.2017.1375989}.

\leavevmode\vadjust pre{\hypertarget{ref-Fecher14}{}}%
Fecher, B., and S. Friesike. 2014. {``Open Science: One Term, Five
Schools of Thought.''} In \emph{Opening Science}, 17--47. Springer,
Cham.

\leavevmode\vadjust pre{\hypertarget{ref-Gilliland16}{}}%
Gilliland, A. J. 2016. {``Setting the Stage.''} In \emph{Introduction to
Metadata}, 3rd ed. Los Angeles, California: Getty Publications.

\leavevmode\vadjust pre{\hypertarget{ref-Hampton15}{}}%
Hampton, S. E., S. S. Anderson, S. C. Bagby, C. Gries, X. Han, E. M.
Hart, M. B. Jones, et al. 2015. {``The Tao of Open Science for
Ecology.''} \emph{Ecosphere} 6 (7): 1--13.
\url{https://doi.org/10.1890/ES14-00402.1}.

\leavevmode\vadjust pre{\hypertarget{ref-Lowenberg21}{}}%
Lowenberg, Daniella, Rachael Lammey, Matthew B Jones, John Chodacki, and
Martin Fenner. 2021. {``Data Citation: Let's Choose Adoption over
Perfection.''} Zenodo. \url{https://doi.org/10.5281/zenodo.4701079}.

\leavevmode\vadjust pre{\hypertarget{ref-Lowndes17}{}}%
Lowndes, J. S. S., B. D. Best, C. Scarborough, J. C. Afflerbach, M. R.
Frazier, C. C. O'Hara, N. Jiang, and B. S. Halpern. 2017. {``Our Path to
Better Science in Less Time Using Open Data Science Tools.''}
\emph{Nature Ecology \& Evolution} 1 (0160): 1--7.
\url{https://doi.org/10.1038/s41559-017-0160}.

\leavevmode\vadjust pre{\hypertarget{ref-Michener97}{}}%
Michener, W. K., J. W. Brunt, J. J. Helly, T. B. Kirchner, and S. G.
Stafford. 1997. {``Nongeospatial Metadata for the Ecological
Sciences.''} \emph{Ecological Applications} 7 (1): 330--42.
https://doi.org/\url{https://doi.org/10.1890/1051-0761(1997)007\%5B0330:NMFTES\%5D2.0.CO;2}.

\leavevmode\vadjust pre{\hypertarget{ref-Wickham14c}{}}%
Wickham, H. 2014. {``Tidy Data.''} \emph{Journal of Statistical
Software} 59 (10): 1--23. \url{https://doi.org/10.18637/jss.v059.i10}.

\leavevmode\vadjust pre{\hypertarget{ref-Wickham19}{}}%
Wickham, H., M. Averick, J. Bryan, W. Chang, L. D'Agostino McGowan, R.
François, G. Grolemund, et al. 2019. {``Welcome to the {tidyverse}.''}
\emph{Journal of Open Source Software} 4 (43): 1686.
\url{https://doi.org/10.21105/joss.01686}.

\leavevmode\vadjust pre{\hypertarget{ref-Wickham17}{}}%
Wickham, H., and G. Grolemund. 2017. \emph{R for Data Science}.
Sebastopol, California: O'Reilly.

\leavevmode\vadjust pre{\hypertarget{ref-Wilkinson16}{}}%
Wilkinson, M. D., M. Dumontier, I. J. Aalbersberg, G. Appleton, M.
Axton, A. Baak, N. Blomberg, et al. 2016. {``The FAIR Guiding Principles
for Scientific Data Management and Stewardship.''} \emph{Scientific
Data} 3 (160018). \url{https://doi.org/10.1038/sdata.2016.18}.

\leavevmode\vadjust pre{\hypertarget{ref-Wilson17}{}}%
Wilson, G., J. Bryan, K. Cranston, J. Kitzes, L. Nederbragt, and T. K.
Teal. 2017. {``Good Enough Practices in Scientific Computing.''}
\emph{PLoS Computational Biology} 13 (6): e1005510.
\url{https://doi.org/10.1371/journal.pcbi.1005510}.

\leavevmode\vadjust pre{\hypertarget{ref-Ziemann16}{}}%
Ziemann, M., Y. Eren, and A. El-Osta. 2016. {``Gene Name Errors Are
Widespread in the Scientific Literature.''} \emph{Genome Biology} 17
(1): 1--3. \url{https://doi.org/10.1186/s13059-016-1044-7}.

\end{CSLReferences}



\end{document}
